\section*{Erd\H{o}s Problem 678: reversing the order of interval LCMs}

\paragraph{Statement and claim.}
Write \(M(n,k)=\operatorname{lcm}(n+1,\ldots,n+k)\).
There are infinitely many triples \(n,m,k\), with \(k\geq3\) and
\(m\geq n+k\), for which \(M(n,k)>M(m,k+1)\).
In fact, for every fixed \(C\geq1\) and all sufficiently large \(k\)
we construct disjoint intervals beginning at \(x<y\) such that
\[
 \operatorname{lcm}(x,\ldots,x+k-1)>
 C\operatorname{lcm}(y,\ldots,y+k).
\]
This reconstructs Cambie's theorem, with a simplified choice of \(y\):
\url{https://arxiv.org/abs/2410.09138},
\url{https://www.erdosproblems.com/678}.
The only asymptotic external input is the prime number theorem, used
to find primes in fixed relative intervals.

\paragraph{The duplication identity.}
Put
\[
 L=\operatorname{lcm}(1,\ldots,k),\qquad
 s=\prod_{p\leq\sqrt{k}}p^{\lfloor\log_p k\rfloor}.
\]
For \(p>\sqrt{k}\), write \(r_p=k\bmod p\), and represent zero
residues by \(p\). Choose positive integers \(x,y\) satisfying
\[
 x\equiv1\pmod s,\quad y\equiv0\pmod s,\qquad
 1\leq x\bmod p\leq p-r_p,\quad
 p-r_p\leq y\bmod p\leq p
 \quad(\sqrt{k}<p\leq k).
                                                               \tag{1}
\]
Let \(D_x\) be the product of the \(k\) terms starting at \(x\),
divided by their LCM, and define \(D_y\) similarly for the \(k+1\)
terms starting at \(y\). Then
\[
 D_y=L D_x.                                                     \tag{2}
\]
Here is a prime-by-prime proof. For \(p\leq k\), set
\(a=\lfloor\log_p k\rfloor\). Each interval contains a multiple
of \(p^a\) and at most one multiple of \(p^{a+1}\). Consequently
the valuation of its duplication factor is the sum of the
valuations truncated at \(a\), minus \(a\).
For \(p\leq\sqrt{k}\), the residues in the \(y\)-interval are
\(0,1,\ldots,k\) modulo \(p^a\), whereas those in the \(x\)-interval
are \(1,\ldots,k\). The difference of the truncated sums is \(a\).
For \(p>\sqrt{k}\), condition (1) makes the \(x\)-interval contain
\(\lfloor k/p\rfloor\) multiples of \(p\), and the \(y\)-interval
one more. For \(p>k\) there is no duplication in either interval.
This proves (2). It follows that
\[
 \frac{\operatorname{lcm}(x,\ldots,x+k-1)}
      {\operatorname{lcm}(y,\ldots,y+k)}
 =\frac L{y+k}\prod_{j=0}^{k-1}\frac{x+j}{y+j}
 \geq \frac L{y+k}\left(\frac xy\right)^k.                       \tag{3}
\]

\paragraph{Choosing two nearby admissible intervals.}
Fix \(\varepsilon=1/(1000C)\). For all sufficiently large \(k\),
choose a prime
\[
 k/2<p_0<(1+\varepsilon)k/2
\]
and three different primes \(p_1,p_2,p_3\) in
\(((1-\varepsilon)k,k)\).
First choose \(y\) among the multiples of \(L/p_0\) in
\[
 \frac L{5C}(1+1/k)<y<\frac L{4C}-k.                            \tag{4}
\]
The number of candidates is
\((1/(20C)+o(1))p_0\). They have distinct residues modulo \(p_0\),
since the interval has length less than \(L\).
The forbidden residues are precisely
\(1,\ldots,2p_0-k-1\), at most \(\varepsilon k\) of them.
There are therefore more candidates than forbidden residues.
All other large-prime residues of \(y\) are zero, which are
allowed in (1).

Put \(P=p_1p_2p_3\), \(g=L/P\), and consider \(x=1+gz\).
All the congruences for \(x\) except those at \(p_1,p_2,p_3\)
are then automatic. At \(p_i\), exactly \(k-p_i\) residues
are forbidden. Define
\[
 h=1+\sum_{i=1}^3(k-p_i)\leq3\varepsilon k+1<\min_i p_i.
\]
In any \(h\) consecutive values of \(z\), its induced residues
modulo each \(p_i\) are distinct. At most \(\sum_i(k-p_i)=h-1\)
of these values are forbidden by any of the three conditions.
Thus every such block contains an admissible \(x\).
The interval
\[
 y-\frac L{5Ck}<x<y                                             \tag{5}
\]
corresponds to a \(z\)-interval of length \(P/(5Ck)\), which is
of order \(k^2\), and hence contains \(h\) consecutive integers.
It therefore contains an admissible \(x\).

Both \(2^{\lfloor\log_2 k\rfloor}>k/2\) and
\(3^{\lfloor\log_3 k\rfloor}>k/3\) divide \(s\), so \(s>k^2/6\).
As \(y-x\equiv-1\pmod s\) and \(y>x\), we have
\(y-x\geq s-1>k\). In particular the intervals are disjoint.
Equations (4)--(5) also give
\[
 x/y>k/(k+1),\qquad y+k<L/(4C).
\]
Substitution in (3) yields a ratio greater than
\[
 4C\left(\frac{k}{k+1}\right)^k>4C/e>C.
\]
Taking \(n=x-1\), \(m=y-1\), and letting \(k\) increase proves
the assertion.

\paragraph{Why \(k\) cannot be held fixed in this conclusion.}
For any interval of \(r\) consecutive positive integers,
its product divided by its LCM is at most \((r-1)!\).
Indeed, for each prime choose a term with maximal valuation.
The valuations of the remaining terms are bounded by those of
their nonzero differences from this term. Their product divides,
prime by prime, \(j!(r-1-j)!\), which divides \((r-1)!\).
If \(m\geq n+k\) and \(M(n,k)>M(m,k+1)\), it follows that
\[
 m^k\geq M(n,k)>M(m,k+1)\geq (m+1)^{k+1}/k!,
\]
so \(m<k!\). There are only finitely many examples for each fixed
\(k\); the infinite family above correctly allows \(k\) to vary.

\paragraph{Completion estimate.}
\textbf{COMPLETION: 100\%}
